\documentclass[11pt,oneside,letterpaper]{article}
\input{preamble.tex}

\title{CORP-PD v4.2 \\ \large Annual Validation Report 2026}
\author{Model Validation \\ \small Large Corporate PD model MRM-0417}
\date{\today}

\begin{document}

\maketitle

\tableofcontents

\clearpage

\section{Executive summary}\label{sec:summary}
This report presents the 2026 annual validation of the CORP-PD v4.2 probability of default
model for large corporate obligors. The model continues to rank-order risk well: the
full-sample AUC remains well above the acceptance threshold (\autoref{tab:discrimination}).
Calibration is weaker. Grade~7 under-predicts observed defaults and the portfolio-level
binomial test rejects at the 1\% level, which leads to Finding F-02 in
\autoref{sec:findings}.

\begin{quirecallout}{Validation opinion · Fit for use with conditions}
The model may remain in use subject to recalibration of the sub-investment-grade tail
(Finding F-02) and remediation of the EBITDA data gap (Finding F-01) by Q2 2027.
\end{quirecallout}

\section{Scope and model overview}\label{sec:scope}
CORP-PD v4.2 assigns one of eight rating grades to large corporate obligors with annual
revenue above USD 500 million. The grade-level PD is the long-run average default rate
mapped through the master scale. The model is used for credit approval, pricing and
regulatory capital, consistent with the IRB requirements in \citep{bcbs_irb_2017}.

The rating score is a logistic function of financial and qualitative factors:
\begin{equation}\label{eq:logit}
  \mathrm{PD}_i = \frac{1}{1 + \exp\!\left(-\beta_0 - \sum_{k=1}^{K} \beta_k x_{ik}\right)}
\end{equation}

\section{Data}\label{sec:data}
\subsection{Sources and lineage}\label{sec:sources}
Obligor financials are taken from the spreading system and ratings from the rating
workflow. The validation sample covers 6,657 obligors rated at the 2025 year end.

\subsection{Observation window}\label{sec:obs-window}
Defaults are observed over the one-year horizon following each rating date between 2019
and 2025. Defaults use the 90-days-past-due or unlikeliness-to-pay definition.

\section{Performance testing}\label{sec:performance}
All tests are executed by the validation suite in \texttt{validation/run\_suite.py}. Each
table below is generated from the latest test run and records the run that produced it.

\subsection{Discrimination}\label{sec:discrimination}
Discrimination is measured with the area under the ROC curve (AUC), the Gini coefficient
and the Kolmogorov--Smirnov statistic, where $\mathrm{Gini} = 2\,\mathrm{AUC} - 1$.

% quire:generated id=tab_discrimination
\input{generated/tab_discrimination.tex}
% quire:end

\subsection{Calibration}\label{sec:calibration}
Calibration was assessed at the rating-grade level by comparing the long-run average
predicted PD with observed one-year default rates. Following \citet{tasche2013}, each
grade was tested with a Jeffreys test at the 5\% level, and the portfolio with the
Hosmer--Lemeshow statistic:
\begin{equation}\label{eq:hosmer-lemeshow}
  \mathrm{HL} = \sum_{g=1}^{G} \frac{(O_g - N_g\bar{p}_g)^2}{N_g\bar{p}_g(1-\bar{p}_g)}
\end{equation}

% quire:generated id=tab_calib_grade
\input{generated/tab_calib_grade.tex}
% quire:end

% quire:generated id=sum_calibration
\input{generated/sum_calibration.tex}
% quire:end

% quire:generated id=fig_reliability
\input{generated/fig_reliability.tex}
% quire:end

\subsection{Stability}\label{sec:stability}
Population stability is measured against the development sample with the population
stability index (PSI) on the score distribution within each sector.

% quire:generated id=tab_psi_sector
\input{generated/tab_psi_sector.tex}
% quire:end

\section{Findings and limitations}\label{sec:findings}
\begin{quirecallout}{Finding F-01 · Severity: Low}
EBITDA is missing for 3.9\% of obligors, against a 2\% tolerance. Missing values are
imputed with sector medians, which flatters leverage-driven scores.
\end{quirecallout}

\begin{quirecallout}{Finding F-02 · Severity: Medium}
Grade~7 under-predicts observed defaults. Recalibrate the grade 6--8 master-scale
boundaries before the 2027 annual review and re-run the Jeffreys test on the recalibrated
scale.
\end{quirecallout}

\bibliographystyle{plainnat}
\bibliography{references}

\appendix

\section{Test suite listing}\label{app:code}
The Jeffreys test used for grade-level calibration:
\begin{lstlisting}[language=Python]
def jeffreys_pvalue(defaults, obligors, pd):
    """One-sided: small p means the PD under-predicts defaults."""
    return betainc(defaults + 0.5, obligors - defaults + 0.5, pd)
\end{lstlisting}

\end{document}
